\documentclass[11pt]{article}
\usepackage[margin=1in]{geometry}
\usepackage{amsmath,amssymb}
\usepackage[hidelinks]{hyperref}
\newcommand{\R}{\mathbb R}
\newcommand{\M}{M}

\title{The classical cube-maximal subcase of two rescaled-space conjectures}
\author{Run \texttt{fa\_banach\_001}, lane 1}
\date{12 August 2026}

\begin{document}
\maketitle

\section*{Status}

\textbf{Literature-implied answer (proper subcase).}
Lerner's 2025 theorem proves both conjectures for Banach function spaces on
$\R^n$ and the usual Hardy--Littlewood maximal operator over cubes.  The source
states the conjectures for a more general quasi-Banach/abstract-basis setting,
which is not covered.  Thus this is neither a new proof nor a full answer to
the source formulations.

\section*{Source questions}

In Z. Nieraeth, \emph{Extrapolation in general quasi-Banach function spaces},
arXiv:2211.03458, Conjectures 2.38 and 2.39 occur together on PDF page 26.
For an $r$-convex, $s$-concave function space $X$, write
\[
 X_{r,s}=\left[\left[(X^r)'\right]^{(s/r)'}\right]'.
\]
Conjecture 2.38 says, for a basis $\mathcal E$ having the $A_1$
self-improvement property and under boundedness of
$M^{\mathcal E}$ on $[(X^r)']^{(s/r)'}$, that
\[
 M^{\mathcal E}:X_{r,s}\to X_{r,s}
 \quad\Longleftrightarrow\quad
 M^{\mathcal E}:X^r\to X^r.
\]
Here the source permits $X$ to be quasi-Banach.  Conjecture 2.39, for Banach
$X$, says that simultaneous boundedness on $X_{r,s}$ and $(X_{r,s})'$ is
equivalent to simultaneous boundedness on $X^r$ and $(X')^{s'}$.

\section*{Later theorem and exact match}

A. K. Lerner, \emph{A note on the maximal operator on Banach function spaces},
J. Geom. Anal. 35 (2025), Article 97, arXiv:2404.15671, explicitly restates
these as Conjectures 1.1 and 1.2.  Lerner assumes
\[
 1\le r<s\le\infty,\qquad X\text{ a Banach function space over }\R^n,
\]
and takes $M$ to be the standard cube Hardy--Littlewood maximal operator.
Section 4, PDF pages 10--11, proves both conjectures.  In particular:
\begin{itemize}
\item under boundedness of $M$ on $[(X^r)']^{(s/r)'}$, boundedness on
  $X_{r,s}$ is equivalent to boundedness on $X^r$;
\item boundedness of $M$ on both $X_{r,s}$ and $(X_{r,s})'$ is equivalent to
  boundedness on both $X^r$ and $(X')^{s'}$.
\end{itemize}
These statements are exactly the source conjectures after specializing
$\Omega=\R^n$, $\mathcal E$ to cubes, $r\ge1$, and $X$ to a Banach function
space.  The identification is explicit in Lerner's introduction, which also
notes that Nieraeth's formulations are more general.

\section*{Why this does not settle the abstract version}

The restriction is substantive.  Lerner's proof uses the local median maximal
operator and a cube-based criterion for dual maximal boundedness.  It also uses
maximal self-improvement, the Fefferman--Stein inequality, $A_p$-regularity,
and reverse-H\"older properties of Euclidean cube weights.  The second
conjecture further passes through a theorem turning two $A_1$-regular
convexifications into $A_q$-regularity.  Nieraeth's sole hypothesis that an
abstract basis has the $A_1$ self-improvement property does not supply these
dual and $A_p$-regularity mechanisms.  The quasi-Banach part of Conjecture 2.38
also lies outside Lerner's Lorentz--Luxembourg and K\"othe-duality argument.

Accordingly, the correct status is a complete affirmative answer in the
canonical Euclidean Banach setting and an unresolved remainder for general
bases and the quasi-Banach range.

\begin{thebibliography}{9}
\bibitem{N}
Z. Nieraeth,
\emph{Extrapolation in general quasi-Banach function spaces},
J. Fourier Anal. Appl. 29 (2023), Article 70; arXiv:2211.03458.

\bibitem{L}
A. K. Lerner,
\emph{A note on the maximal operator on Banach function spaces},
J. Geom. Anal. 35 (2025), Article 97; arXiv:2404.15671.
\end{thebibliography}

\end{document}
